\section*{Capítulo \ref{ch:b2:atomic-orbitals} --- Orbitais atômicos: formas, energias e tamanhos}

\begin{solution}{exo:b2:atomic-orbitals:1}
3p: $n = 3$, $l = 1$; um \omterm{def:b2:atomic-orbitals:node}{nó radial} ($n - l - 1$) e um \omterm{def:b2:atomic-orbitals:node}{nó angular}. 4d: $n = 4$,
$l = 2$; um \omterm{def:b2:atomic-orbitals:node}{nó radial} e dois \omterm{def:b2:atomic-orbitals:node}{nós angulares}.
\end{solution}

\begin{solution}{exo:b2:atomic-orbitals:2}
$P_{2p} \propto r^4\mathrm e^{-r}$ ($r$ em $a_0$): $d\ln P/dr = 4/r - 1 = 0$ para $r = 4$:
$4a_0 = \qty{212}{pm}$.
\end{solution}

\begin{solution}{exo:b2:atomic-orbitals:3}
$1 - \mathrm e^{-2}(1 + 2 + 2) = 1 - 5\mathrm e^{-2} = 0.323$.
\end{solution}

\begin{solution}{exo:b2:atomic-orbitals:4}
Quatro lobos que apontam entre os eixos $x$ e $y$, no plano $xy$, de sinais
alternados (positivos onde $xy > 0$); os planos nodais são $xz$ e $yz$. \omterm{def:b2:atomic-orbitals:node}{Nós
radiais}: $n - l - 1 = 0$.
\end{solution}

\begin{solution}{exo:b2:atomic-orbitals:5}
Oxigênio, $(1\mathrm s)^2(2\mathrm s, 2\mathrm p)^6$: $\sigma = 5 \times 0.35 + 2 \times 0.85
= 3.45$, $Z^* = 4.55$; raio $4/4.55 = 0.88a_0$, \qty{47}{pm}.
\end{solution}

\begin{solution}{exo:b2:atomic-orbitals:6}
A ionização retira os elétrons menos ligados. Os elétrons 4s
(\qty{-14}{eV}) estão muito menos ligados que os 3d (\qty{-59}{eV}): os dois
elétrons 4s saem primeiro, depois um elétron 3d, o que dá $[\ce{Ar}]\,3\mathrm d^5$.
\end{solution}

\begin{solution}{exo:b2:atomic-orbitals:7}
Lítio, 2s: $Z^* = 3 - 2 \times 0.85 = 1.30$, $I = 13.6 \times 1.30^2/4 =
\qty{5.75}{eV}$ (medido 5.39, 7~\% acima). Sódio, 3s: $Z^* = 11 - 8 \times 0.85 -
2 \times 1.00 = 2.20$, $I = 13.6 \times 2.20^2/9 = \qty{7.3}{eV}$ (medido 5.14): as
regras, um modelo grosseiro, superestimam a ligação do elétron externo, e
cada vez mais descendo o grupo.
\end{solution}

\begin{solution}{exo:b2:atomic-orbitals:8}
Sódio: $Z^* = 2.20$, raio $9/2.20 = 4.1a_0$ (\qty{216}{pm}). Cloro, 3p:
$\sigma = 6 \times 0.35 + 8 \times 0.85 + 2 \times 1.00 = 10.90$, $Z^* = 6.10$, raio
$9/6.10 = 1.5a_0$ (\qty{78}{pm}). Mesma camada, mas o elétron externo do cloro
sente uma carga quase três vezes maior e é puxado para dentro.
\end{solution}

\begin{solution}{exo:b2:atomic-orbitals:9}
Ambos têm a \omterm{def:b2:atomic-orbitals:radial-angular}{parte angular} $1/\sqrt{4\pi}$; logo a integral se reduz a
\[
  \int_0^\infty R_{1s}R_{2s}r^2\,dr \propto \int_0^\infty (2 - r)r^2\mathrm e^{-3r/2}\,dr =
  2 \times \frac{2!}{(3/2)^3} - \frac{3!}{(3/2)^4} = \frac{32}{27} - \frac{32}{27} = 0 .
\]
\end{solution}

\begin{solution}{exo:b2:atomic-orbitals:10}
$\int_0^\infty N^2r^2\mathrm e^{-r}r^2\,dr = N^2 \times 4! = 1$: $N = 1/\sqrt{24} =
1/(2\sqrt6)$, como na tabela.
\end{solution}

\begin{solution}{exo:b2:atomic-orbitals:11}
$\langle r\rangle = 1.5a_0$, \omterm{def:b2:atomic-orbitals:radial-distribution}{raio mais provável} $a_0$. $P(r) = 4r^2\mathrm e^{-2r}$ sobe
abruptamente a partir de zero e cai lentamente: a longa cauda em $r$ grande
pesa mais na média do que na posição do máximo.
\end{solution}

\begin{solution}{exo:b2:atomic-orbitals:12}
$Y_{p_z} \propto z/r = \cos\theta$ se anula para $\theta = \pi/2$: o plano $xy$.
$\psi$ tem o sinal de $z$, positivo acima, negativo abaixo. No plano, a
\omterm{def:b2:atomic-orbitals:wavefunction}{densidade de probabilidade} é nula.
\end{solution}

\begin{solution}{pb:b2:atomic-orbitals:1}
\textbf{1.} $\int|\psi|^2\,dV = \frac{4\pi}{\pi a_0^3}\int_0^\infty r^2\mathrm e^{-2r/a_0}\,dr =
\frac{4}{a_0^3} \times \frac{a_0^3}{4} = 1$.
\textbf{2.} $P(r) = 4\pi r^2|\psi|^2 = (4/a_0^3)\,r^2\mathrm e^{-2r/a_0}$.
\textbf{3.} $dP/dr = 0$: $2/r - 2/a_0 = 0$, $r = a_0$.
\textbf{4.} $\langle r\rangle = (4/a_0^3)\int_0^\infty r^3\mathrm e^{-2r/a_0}\,dr = (4/a_0^3)
\times 6(a_0/2)^4 = 1.5a_0$.
\textbf{5.} A distribuição é assimétrica, com uma longa cauda em $r$ grande.
\textbf{6.} $|\psi|^2$ é uma densidade por unidade de volume, máxima no núcleo;
$P(r)$ conta o volume da casca, $4\pi r^2\,dr$, que se anula em $r = 0$.
\textbf{7.} \qty{52.9}{pm}.
\textbf{8.} Com $x = r/a_0$: $\int_{x_0}^\infty 4x^2\mathrm e^{-2x}\,dx$; integrando por
partes duas vezes, $= \mathrm e^{-2x_0}(2x_0^2 + 2x_0 + 1)$.
\textbf{9.} $5\mathrm e^{-2} = 0.677$.
\textbf{10.} $13\mathrm e^{-4} = 0.238$.
\textbf{11.} $\mathrm e^{-2x}(1 + 2x + 2x^2) = 0.10$ para $x \approx 2.66$: \qty{141}{pm}.
\textbf{12.} A densidade nunca se anula, a nenhuma distância: não há borda. Mas
90~\% da probabilidade está a menos de \qty{141}{pm}: um tamanho prático.
\textbf{13.} 2s: um \omterm{def:b2:atomic-orbitals:node}{nó radial} (esfera de raio $2a_0$), nenhum \omterm{def:b2:atomic-orbitals:node}{nó angular}. 2p: um
\omterm{def:b2:atomic-orbitals:node}{nó angular} (um plano), nenhum \omterm{def:b2:atomic-orbitals:node}{nó radial}.
\textbf{14.} Em $r = 2a_0$.
\textbf{15.} $4a_0$ (exercício 2).
\textbf{16.} $d\ln P/dr = 2/r - 2/(2 - r) - 1 = 0$ dá $r^2 - 6r + 4 = 0$, $r = 3 \pm \sqrt5$:
$0.76a_0$ (pequeno máximo interno) e $5.24a_0$ (máximo principal).
\textbf{17.} O 2s, por seu máximo interno. Em um átomo polieletrônico, um
elétron 2s penetra na camada 1s, é menos blindado que um elétron 2p e fica
mais baixo em energia.
\textbf{18.} Todos os raios são divididos por $Z = 6$: 2p em $0.67a_0$ (\qty{35}{pm});
máximos do 2s em $0.13a_0$ e $0.87a_0$.
\textbf{19.} C 3.25, N 3.90, O 4.55.
\textbf{20.} $\varepsilon = -13.6\,Z^{*2}/4$: C \qty{-35.9}{eV}, N \qty{-51.7}{eV}, O
\qty{-70.4}{eV}.
\textbf{21.} Como no capítulo: \qty{11.5}{eV} (medido 11.26).
\textbf{22.} N: atom $5 \times (-13.6 \times 3.90^2/4) = \qty{-258.6}{eV}$; \ce{N+}, $Z^* =
4.25$, $4 \times (-13.6 \times 4.25^2/4) = \qty{-245.7}{eV}$; $I = \qty{12.9}{eV}$.
\textbf{23.} O: atom $6 \times (-13.6 \times 4.55^2/4) = \qty{-422.3}{eV}$; \ce{O+}, $Z^* =
4.90$, $5 \times (-13.6 \times 4.90^2/4) = \qty{-408.2}{eV}$; $I = \qty{14.2}{eV}$. Modelo
11.5, 12.9, 14.2; medidos 11.26, 14.53, 13.62.
\textbf{24.} A energia medida do oxigênio é menor que a do nitrogênio: no
oxigênio, o quarto elétron 2p precisa compartilhar um orbital com outro, e a
repulsão entre eles o torna mais fácil de retirar, enquanto os três elétrons
p desemparelhados do nitrogênio formam uma subcamada semipreenchida estável. As
\omterm{def:b2:atomic-orbitals:slater}{regras de Slater} tratam todos os elétrons de um grupo da mesma forma e nada
sabem do emparelhamento.
\textbf{25.} O elétron 1s do hidrogênio está além de $2a_0$ com probabilidade
$\boldsymbol{13\mathrm e^{-4} \approx 0.238}$.
\end{solution}
